\newpage
\section{第18章：案例研究——命令式对象}

\begin{taplauthorsolutionentrybox}
\phantomsection
\label{sol:author:18.6.1}
\textbf{18.6.1 解答}

\begin{lstlisting}
DecCounter =
  {get:Unit->Nat, inc:Unit->Unit, reset:Unit->Unit,
   dec:Unit->Unit};

decCounterClass =
  lambda r:CounterRep.
    let super = resetCounterClass r in
      {get = super.get,
       inc = super.inc,
       reset = super.reset,
       dec = lambda _:Unit. r.x := pred(!(r.x))};
\end{lstlisting}

\taplsolutionbacklink{ex:18.6.1}{返回练习18.6.1}
\end{taplauthorsolutionentrybox}

\begin{taplauthorsolutionentrybox}
\phantomsection
\label{sol:author:18.7.1}
\textbf{18.7.1 解答}

\begin{lstlisting}
BackupCounter2 =
  {get:Unit->Nat, inc:Unit->Unit,
   reset:Unit->Unit, backup:Unit->Unit,
   reset2:Unit->Unit, backup2:Unit->Unit};

BackupCounterRep2 = {x:Ref Nat, b:Ref Nat, b2:Ref Nat};

backupCounterClass2 =
  lambda r:BackupCounterRep2.
    let super = backupCounterClass r in
      {get = super.get,
       inc = super.inc,
       reset = super.reset,
       backup = super.backup,
       reset2 = lambda _:Unit. r.x := !(r.b2),
       backup2 = lambda _:Unit. r.b2 := !(r.x)};
\end{lstlisting}

\taplsolutionbacklink{ex:18.7.1}{返回练习18.7.1}
\end{taplauthorsolutionentrybox}

\begin{taplauthorsolutionentrybox}
\phantomsection
\label{sol:author:18.11.1}
\textbf{18.11.1 解答}

先让 \texttt{get} 和 \texttt{set} 都递增访问计数：

\begin{lstlisting}
instrCounterClass =
  lambda r:InstrCounterRep.
  lambda self:Unit->InstrCounter.
  lambda _:Unit.
    let super = setCounterClass r self unit in
      {get = lambda _:Unit.
               (r.a := succ(!(r.a)); super.get unit),
       set = lambda i:Nat.
               (r.a := succ(!(r.a)); super.set i),
       inc = super.inc,
       accesses = lambda _:Unit. !(r.a)};
\end{lstlisting}

加入 \texttt{reset} 的子类如下：

\begin{lstlisting}
ResetInstrCounter =
  {get:Unit->Nat, set:Nat->Unit, inc:Unit->Unit,
   accesses:Unit->Nat, reset:Unit->Unit};

resetInstrCounterClass =
  lambda r:InstrCounterRep.
  lambda self:Unit->ResetInstrCounter.
  lambda _:Unit.
    let super = instrCounterClass r self unit in
      {get = super.get,
       set = super.set,
       inc = super.inc,
       accesses = super.accesses,
       reset = lambda _:Unit. r.x := 0};
\end{lstlisting}

最后加入备份功能，并给出对象生成器：

\begin{lstlisting}
BackupInstrCounter =
  {get:Unit->Nat, set:Nat->Unit, inc:Unit->Unit,
   accesses:Unit->Nat, backup:Unit->Unit, reset:Unit->Unit};

BackupInstrCounterRep =
  {x:Ref Nat, a:Ref Nat, b:Ref Nat};

backupInstrCounterClass =
  lambda r:BackupInstrCounterRep.
  lambda self:Unit->BackupInstrCounter.
  lambda _:Unit.
    let super = resetInstrCounterClass r self unit in
      {get = super.get,
       set = super.set,
       inc = super.inc,
       accesses = super.accesses,
       reset = lambda _:Unit. r.x := !(r.b),
       backup = lambda _:Unit. r.b := !(r.x)};

newBackupInstrCounter =
  lambda _:Unit.
    let r = {x=ref 1, a=ref 0, b=ref 0} in
      fix (backupInstrCounterClass r) unit;
\end{lstlisting}

\taplsolutionbacklink{ex:18.11.1}{返回练习18.11.1}
\end{taplauthorsolutionentrybox}

\begin{taplauthorsolutionentrybox}
\phantomsection
\label{sol:author:18.13.1}
\textbf{18.13.1 解答}

可以用引用单元检验对象身份。先在对象的内部表示中加入类型为
\texttt{Ref (Ref Nat)} 的身份字段：

\begin{lstlisting}
IdCounterRep = {x:Ref Nat, id:Ref (Ref Nat)};

IdCounter =
  {get:Unit->Nat, inc:Unit->Unit, id:Unit->Ref Nat};

idCounterClass =
  lambda r:IdCounterRep.
    {get = lambda _:Unit. !(r.x),
     inc = lambda _:Unit. r.x := succ(!(r.x)),
     id = lambda _:Unit. !(r.id)};
\end{lstlisting}

\texttt{id} 方法只返回身份字段中保存的引用。\texttt{sameObject} 接受两个
具有 \texttt{id} 方法的对象，并检查它们返回的引用是否相同：

\begin{lstlisting}
sameObject =
  lambda a:{id:Unit->Ref Nat}.
  lambda b:{id:Unit->Ref Nat}.
    ((b.id unit) := 1;
     (a.id unit) := 0;
     iszero (!(b.id unit)));
\end{lstlisting}

这里利用别名来检验两条引用是否指向同一单元：先确保第二条引用保存非零值，
再向第一条写入零，最后检查第二条是否也变成零。

\taplsolutionbacklink{ex:18.13.1}{返回练习18.13.1}
\end{taplauthorsolutionentrybox}
